\documentclass[10pt,twocolumn]{article}
\usepackage[a4paper,margin=15mm]{geometry}
\usepackage{amsmath,amssymb}
\usepackage{amsthm}
\theoremstyle{definition}
\newtheorem{theorem}{Theorem}
\newtheorem{proposition}{Proposition}
\newtheorem{corollary}{Corollary}
\newtheorem{remark}{Remark}
\newtheorem{example}{Example}
\newtheorem{definition}{Definition}
\usepackage{graphicx}
\usepackage{booktabs,array,calc}
\usepackage[font=small,labelfont=bf]{caption}
\usepackage[colorlinks=true]{hyperref}
\usepackage{etoolbox}
\providecommand{\tightlist}{\setlength{\itemsep}{0pt}\setlength{\parskip}{0pt}}
\setcounter{secnumdepth}{-1}
\emergencystretch=3em
\allowdisplaybreaks
\AtBeginDocument{%
  \BeforeBeginEnvironment{equation}{\small}%
  \BeforeBeginEnvironment{equation*}{\small}%
  \BeforeBeginEnvironment{align}{\small}%
  \BeforeBeginEnvironment{align*}{\small}%
  \BeforeBeginEnvironment{gather}{\small}%
  \BeforeBeginEnvironment{gather*}{\small}%
}

\begin{document}
\title{Finite Nonviability Certificates under Delayed Information}
\author{Amin Abaee\\[0.35em]
{\small Independent Researcher}\\[0.55em]
{\small\href{https://orcid.org/0000-0002-0019-1842}{ORCID: 0000-0002-0019-1842}}\\[0.3em]
{\small\href{mailto:amin\_abaee@ut.ac.ir}{amin\_abaee@ut.ac.ir}}}
\date{September 23, 2026}
\maketitle

\begin{abstract}

We establish a continuous-to-finite theorem for robust nonviability under
exogenous partial observation. An outer moment approximation of all measurable
information-adapted controls, combined with an inner adversarial approximation,
yields finite linear programs whose optimal values sandwich the continuous-time
safety value between explicit error estimates. Every nonviable initial
information state consequently admits a finite, independently verifiable
obstruction certificate, and finite termination holds for every nonviable
instance in the stated class. The certificates account for post-observation
recourse: they can detect failure even when every compatible state is viable
under full information and a common control keeps every branch safe throughout
the blind window. Certificate sparsity is governed by information--time rank,
and a matching construction shows that instantaneous input dimension cannot
bound the number of necessary witnesses. A three-branch delayed-observation
instance is carried exactly --- delay threshold, Farkas certificate, and mesh
refinement --- and the finite programs are solved by a standard solver, with
optimality certified in exact rational arithmetic by independent primal and
dual witnesses.

\end{abstract}

\subsection{1. Introduction}\label{introduction}

Viability theory under full observation has a mature sufficiency theory; under
incomplete observation, the complementary question is when failure can be
\emph{certified} --- finite, checkable witnesses that no observation-based
policy is viable. The companion paper
\mbox{(Abaee, 2026)} develops that calculus of obstruction certificates for
discrete and continuous systems and names the failure mode its discrete
certificates do not reach: a system can remain safe throughout the blind
window and every compatible state can be viable under full information, yet
the branches that could not be distinguished may have accumulated obligations
that no post-observation control can discharge. This paper supplies the
continuous-time instrument for exactly that mode: finite linear programs whose
infeasibility certifies failure of \emph{every measurable
information-adapted controller}, not merely of sampled or parameterized
controllers, with an explicit approximation-error bound, finite termination
whenever nonviability holds, and a sparsity law governed by
information--time rank.

The distinction that carries the paper is that the certificate may use safety
obligations \emph{after} information arrives: it prices the controller's best
post-observation recourse. It can therefore prove nonviability when every
compatible state is individually viable under full information, when a common
control keeps every branch safe throughout the blind interval, and when
instantaneous admissibility and tangency tests find no obstruction.

Several ingredients are standard and are not advertised as new: robust safety
as a family of linear inequalities, nonanticipativity constraints, an
infinite-dimensional separating functional, Farkas' lemma and Helly's theorem,
and monotonicity under observation refinement. The theorem-driven content is
that one construction establishes, simultaneously: (i) the correct abstraction
direction for nonviability --- outer approximation of controller
possibilities, inner approximation of adversarial tests; (ii) coverage of all
measurable adapted controls, with no hidden zero-order-hold restriction;
(iii) a computable two-sided error bound between a finite linear program and
the continuous-time safety value; (iv) finite termination for every nonviable
instance in the stated class; (v) a finite, policy-independent adversarial
test set with numerical verification margins; and (vi) a sharp
information--time sparsity law, including a construction showing that failure
of any input-dimension-only bound.

Section 2 fixes the model and quantifier order. Sections 3--4 give the exact
robust rows and the finite atomic certificate. Section 5 states the
continuous-to-finite theorem. Section 6 carries a three-branch instance
exactly and reports a solver campaign in which the finite programs are solved
by HiGHS \mbox{(Huangfu and Hall, 2018)} and optimality is certified in exact
rational arithmetic by independent witnesses. Section 7 records complexity and
the certificate-verification rules. Section 8 establishes the sharp limitation
that bounds by input dimension fail for blind control functions. Section 9
derives the refinement-erosion identity and the oracle-recourse extension.
Section 10 delimits scope.

\subsection{2. Model, information, and quantifiers}\label{model}

\subsubsection{2.1 Dynamics}\label{dynamics}

Let the hidden mode be
\(\theta \in \Theta := \{1,\dots,M\}\). Conditional on \(\theta\), the
continuous state satisfies
\[
\dot x(t) = A_\theta(t)x(t) + g_\theta(t) + B_\theta(t)u(t) + E_\theta(t)d(t),
\qquad t \in [0,T],
\]
under the standing assumptions: \(A_\theta, g_\theta, B_\theta, E_\theta\)
bounded and piecewise continuous, with bounded-variation bounds on
\(B_\theta, E_\theta\); the initial set in mode \(\theta\) a nonempty compact
polytope \(P_\theta \subset \mathbb{R}^n\); the input and disturbance sets
nonempty compact polytopes \(U = \{v : Fv \le f\} \subset \mathbb{R}^m\) and
\(D_\theta \subset \mathbb{R}^p\); safety the closed polyhedron
\(\mathcal{V} = \{x : c_i^\top x \le b_i,\ i = 1,\dots,s\}\); and the initial
information state
\(B_0 = \bigcup_{\theta} \{\theta\} \times P_\theta\) initially safe (a
time-zero certificate suffices otherwise). The dynamics admit unique
Carath\'eodory solutions for every admissible measurable signal pair.

\subsubsection{2.2 Information structure}\label{information}

At known times \(0 = \tau_0 < \tau_1 < \dots < \tau_R = T\) the controller
receives \emph{exogenous} information about \(\theta\): refining history
partitions
\(\mathcal{P}_0 \preceq \mathcal{P}_1 \preceq \dots \preceq \mathcal{P}_{R-1}\)
of \(\Theta\), the observation history during
\(I_r = [\tau_r, \tau_{r+1})\) identifying the cell \(C_r(\theta) \in
\mathcal{P}_r\). These are history partitions, not necessarily partitions
generated by the latest observation alone; nonrefining raw observations are
replaced by their history refinement. Observations do not depend on
\(x, u, d\): active sensing and endogenous observation branching are excluded
from the exact theorem (Section~9.2 gives a sound extension).

\subsubsection{2.3 Policy class and quantifier order}\label{policy}

Policies are deterministic, measurable, nonanticipative, and adapted to the
observation history; controls may vary arbitrarily measurably between
observations. Because each partition cell has a unique ancestor history, every
such policy is represented exactly by signals
\(u_{r,C} \in L^\infty(I_r; U)\), \(C \in \mathcal{P}_r\), with
\(\nu^\pi_\theta(t) = u_{r,C_r(\theta)}(t)\) on \(I_r\). The robust quantifier
order is
\[
\exists \pi \in \Pi_{\mathcal{I}}\ \forall \theta\ \forall x_0 \in P_\theta\ 
\forall d \in L^\infty([0,T]; D_\theta)\ \forall t \in [0,T],
\]
the adversary selecting mode, initial state, and measurable disturbance with
knowledge of the announced policy; randomized policies are not included. The
information-constrained viability kernel is
\[
\mathcal{K}_{\mathcal{I}}^T = \Bigl\{ B_0 : \exists \pi\ \forall \theta, x_0,
d, t \quad x^{\pi}_{\theta, x_0, d}(t) \in \mathcal{V} \Bigr\},
\]
and a finite-horizon nonviability certificate also excludes infinite-horizon
viability.

\subsubsection{2.4 Relation to full-information viability}\label{fullinfo}

Let \(K^T_{\mathrm{full}}\) be the physical-state viability kernel when
\((\theta, x)\) is observed continuously, and define its belief lifting by
\(\mathrm{Lift}(K^T_{\mathrm{full}}) = \{ B : \varnothing \ne B \subseteq
K^T_{\mathrm{full}} \}\). Then
\[
\mathcal{K}_{\mathcal{I}}^T \subseteq \mathrm{Lift}(K^T_{\mathrm{full}}),
\]
since a full-information controller can emulate any information-constrained
policy. The inclusion is strict in general: the instance of Section~6 has
every two-state subbelief viable while the three-state belief is not.

\subsection{3. Adjoint safety rows and the safety value}\label{rows}

Let \(X_\theta(t,s)\) denote the state-transition matrix. For a label
\(a = (\theta, i, t) \in \mathscr{S} := \Theta \times \{1,\dots,s\} \times
[0,T]\) define the input kernel and margin
\[
k_a(\sigma) = \mathbf{1}_{\{\sigma \le t\}} B_\theta(\sigma)^\top
X_\theta(t,\sigma)^\top c_i,
\]
\[
\beta_a = b_i - h_{P_\theta}\!\bigl(X_\theta(t,0)^\top c_i\bigr)
- \int_0^t c_i^\top X_\theta(t,\sigma) g_\theta(\sigma)\, d\sigma
- \int_0^t h_{D_\theta}\!\bigl(E_\theta(\sigma)^\top
X_\theta(t,\sigma)^\top c_i\bigr)\, d\sigma,
\]
where \(h_Q(v) = \sup_{q \in Q} v^\top q\). For a policy \(\pi\) set
\(\mathcal{F}_a(\pi) = \int_0^T k_a(\sigma)^\top u^\pi_\theta(\sigma)\,
d\sigma - \beta_a\).

\begin{proposition}[robust rows]\label{prop:rows}
\(\pi\) is robustly safe on \([0,T]\) if and only if
\(\mathcal{F}_a(\pi) \le 0\) for every \(a \in \mathscr{S}\). This follows
from the linear solution formula; initial-state and disturbance maximizations
separate from the control term because the control signal in a mode does not
depend on \(x_0\) or \(d\). For polytopic \(D_\theta\), a maximizing
disturbance can be chosen measurably by selecting a maximizing vertex at each
time.
\end{proposition}

Define the minimax safety violation
\(J_{\mathcal{I}}(T) = \min_{\pi \in \Pi_{\mathcal{I}}} \max_{a \in
\mathscr{S}} \mathcal{F}_a(\pi)\).

\begin{proposition}[value and nonviability]\label{prop:value}
The minimum exists: the policy-signal space is weak-* compact and each
\(\mathcal{F}_a\) is weak-* continuous. Consequently
\[
B_0 \notin \mathcal{K}_{\mathcal{I}}^T
\iff
J_{\mathcal{I}}(T) > 0.
\]
The strict positive margin is a consequence of compactness and closed safety
constraints, not an extra nonviability assumption.
\end{proposition}

\subsection{4. The finite obstruction certificate}\label{certificate}

Define \(\phi_U(g) := \min_{v \in U} g^\top v = -h_U(-g)\). Choose finitely
many labels \(a_\ell = (\theta_\ell, i_\ell, t_\ell)\), \(\ell = 1,\dots,q\),
and weights \(\lambda_\ell \ge 0\), \(\sum_\ell \lambda_\ell = 1\). For each
active information cell \(C \in \mathcal{P}_r\), pool the adjoint kernels
\(g_{r,C}(\sigma) = \sum_{\ell : \theta_\ell \in C} \lambda_\ell
k_{a_\ell}(\sigma)\). The certificate margin is
\[
\Gamma_{\mathcal{I}}(\lambda, a)
= \sum_{r=0}^{R-1} \int_{I_r} \sum_{C \in \mathcal{P}_r}
\phi_U\!\bigl(g_{r,C}(\sigma)\bigr)\, d\sigma
- \sum_{\ell=1}^q \lambda_\ell \beta_{a_\ell}.
\tag{1}
\]

\begin{remark}[semantics]\label{rem:semantics}
If \(\Gamma_{\mathcal{I}}(\lambda, a) > 0\), then no admissible
information-adapted policy is safe: for every policy,
\(\sum_\ell \lambda_\ell \mathcal{F}_{a_\ell}(\pi) \ge
\Gamma_{\mathcal{I}}(\lambda, a)\), so at least one stored label satisfies
\(\mathcal{F}_{a_\ell}(\pi) \ge \Gamma_{\mathcal{I}}(\lambda, a) > 0\) ---
an admissible realization violates facet \(i_\ell\) at time \(t_\ell\). The
violating label may depend on the policy; the finite list and its weights do
not. The certificate is not an adversarial tree: it is a finite,
policy-independent set of adversarial tests combined by an adjoint identity.
\end{remark}

\subsection{5. The continuous-to-finite theorem}\label{bridge}

The exact dual expression (1) alone would be close to standard
infinite-dimensional convex duality. The consequential result is a sound
finite construction covering unrestricted measurable controls.

\subsubsection{5.1 Outer control moments}\label{outer}

Choose a control-time mesh \(\mathcal{H} = \{J_1,\dots,J_H\}\) containing
every sensing time, so each \(J_j\) lies within one information stage.
Introduce one vector \(v_{j,C} \in U\) for every cell \(C\) active during
\(J_j\), and let \(Q = \sum_j |\mathcal{P}_{r(j)}|\) be the number of
information--time blocks. For a label \(a = (\theta, i, t)\) define the
finite row
\(A_{a,j,C} = \mathbf{1}_{\{\theta \in C\}} \int_{J_j} k_a(\sigma)\,
d\sigma\), the block average
\(\bar{k}_{a,\mathcal{H}}(\sigma) = \frac{1}{|J_j|} \int_{J_j} k_a(s)\, ds\)
on \(J_j\), and the moment error
\[
e_{a,\mathcal{H}} = \int_0^T h_U\!\bigl(\bar{k}_{a,\mathcal{H}}(\sigma)
- k_a(\sigma)\bigr)\, d\sigma,
\tag{2}
\]
for which any verified upper bound \(\bar{e}_a \ge e_{a,\mathcal{H}}\) may be
used. For every measurable policy, its block averages
\(v_{j,C} = \frac{1}{|J_j|} \int_{J_j} u_{r(j),C}(\sigma)\, d\sigma\) belong
to \(U\), and
\[
A_a v - \int_0^T k_a(\sigma)^\top u^\pi_\theta(\sigma)\, d\sigma
\le \bar{e}_a.
\tag{3}
\]
Inequality (3) is the crucial orientation: the finite moments yield a
\emph{necessary relaxation} of continuous-time safety. They do not restrict
the real controller to held inputs. (For the sampled-data subclass whose
policies are held constant on the blocks, the same polyhedron is an exact
reduction rather than a relaxation --- the held class is the policy class, so
only inter-sample times are relaxed, which is sound.)

\subsubsection{5.2 Inner adversarial approximation}\label{inner}

One may use a right-hand side \(\beta_a^{+} \ge \beta_a\) obtained from an
explicit admissible initial state and disturbance signal rather than exact
worst-case optimization. If \(0 \le \beta_a^{+} - \beta_a \le \delta_\beta\)
uniformly over the labels used, the approximation is safe for nonviability:
it weakens the adversary and enlarges the feasible controller set. For
example, with \(w_a(\sigma) = \mathbf{1}_{\{\sigma \le t\}}
E_\theta(\sigma)^\top X_\theta(t,\sigma)^\top c_i\), on a disturbance mesh
choose \(d_{a,J} \in \arg\max_{d \in D_\theta} (\int_J w_a)^{\top} d\),
yielding an explicit piecewise-constant admissible disturbance; if
\(D_\theta\) has radius \(R_D\) about a fixed center and
\(\mathrm{TV}(w_a) \le V_D\), then
\[
0 \le \beta_a^{+} - \beta_a \le R_D V_D h_d,
\tag{4}
\]
apart from any separately accounted initial-state or numerical error.

\subsubsection{5.3 The finite program}\label{finlp}

Let \(G \subset [0,T]\) be a finite time grid including every mode and safety
facet at every grid time. Solve
\[
\rho_{\mathcal{H},G} = \min_{v, \rho}\ \rho
\quad \text{s.t.} \quad
A_a v - \beta_a^{+} - \bar{e}_a \le \rho \ \ (a \in \Theta \times \{1,\dots,s\} \times G),
\qquad v_{j,C} \in U.
\tag{5}
\]

\begin{theorem}[continuous-to-finite]\label{thm:bridge}
Assume: uniformly over admissible policies, modes, and facets,
\(|\mathcal{F}_{\theta,i,t}(\pi) - \mathcal{F}_{\theta,i,t'}(\pi)| \le
L|t - t'|\) (6); the input kernels satisfy \(\mathrm{TV}(k_a) \le V_U\) (7);
\(U\) has radius \(R_U\) about some fixed center; and the time grid has
covering radius at most \(h_t\), the control mesh maximum length \(h_u\), and
the adversarial approximation satisfies the bound above. Then:
\end{theorem}

\emph{(i) Certified value sandwich.}
\[
\rho_{\mathcal{H},G} \le J_{\mathcal{I}}(T) \le \rho_{\mathcal{H},G}
+ R_U V_U h_u + \delta_\beta + L h_t.
\tag{8}
\]
With (4), \(J_{\mathcal{I}}(T) - \rho_{\mathcal{H},G} \le R_U V_U h_u + R_D
V_D h_d + L h_t\). Verified numerical row widths are added to the right side.

\emph{(ii) Sound finite certificates.} If \(\rho_{\mathcal{H},G} > 0\), the
program's dual supplies a finite certificate of nonviability of the original
continuous-time system, covering every measurable information-adapted policy.

\emph{(iii) Completeness under refinement.} If
\(B_0 \notin \mathcal{K}_{\mathcal{I}}^T\), then every sufficiently fine
sequence of the above meshes, with numerical errors tending to zero,
eventually produces a positive finite certificate. Nonviability is thus
semidecidable in this class when the coefficients and error bounds are
effectively computable.

\emph{(iv) Exact atomic dual representation.}
\[
J_{\mathcal{I}}(T) = \sup_{q < \infty,\ a_1,\dots,a_q,\
\lambda \in \Delta_q} \Gamma_{\mathcal{I}}(\lambda, a).
\tag{10}
\]
In particular, nonviability is equivalent to the existence of a positive
finite atomic certificate.

\emph{(v) Sparse witnesses.} Let \(r = \mathrm{rank}\, A\) for the finite
safety-row matrix in (5). Whenever \(\rho_{\mathcal{H},G} > 0\), a positive
certificate exists using at most
\(r + 1 \le mQ + 1\) safety labels, and it can retain any prescribed margin
strictly below \(\rho_{\mathcal{H},G}\). If the lifted continuous-time
kernels span a finite-dimensional space of dimension \(r_{\mathcal{I}}\), a
continuous-time obstruction has at most \(r_{\mathcal{I}} + 1\) labels. The
relevant dimension is information--time rank, not \(m\).

The complete proof is in Supplementary S1. The asymmetry at its core is
worth stating plainly: averaging arbitrary controls proves the lower bound;
implementing piecewise-constant controls proves the upper bound. Confusing
these two directions is what makes na\"ive infeasibility-by-discretization
arguments unsound.

\subsection{6. A three-branch delayed-observation instance}\label{instance}

The instance is analytically checkable, polyhedral, and genuinely three-way:
every pair of compatible states is viable, but the triple is not.

\subsubsection{6.1 Plant and exact delay threshold}\label{plant}

Consider the planar double integrator \(\dot p = v\), \(\dot v = u\), with
unit vectors \(n_1 = (1,0)\), \(n_2 = (-\tfrac35, \tfrac45)\),
\(n_3 = (-\tfrac35, -\tfrac45)\), the hexagonal input set
\(U = \mathrm{conv}\{\pm n_1, \pm n_2, \pm n_3\}\), equivalently
\(\{\nu : |u_2| \le \tfrac45,\ |2u_1 + u_2| \le 2,\ |2u_1 - u_2| \le 2\}\),
safety \(n_i^\top p \le 2\) (\(i = 1,2,3\)) and \(|v_1|, |v_2| \le \tfrac65\),
three possible initial states \(p_j(0) = \tfrac{34}{25} n_j\),
\(v_j(0) = n_j\), no disturbance, and the mode unobserved until time \(\tau\),
when it is revealed exactly. Every initial state is strictly inside the
physical safety set.

If \(j\) is known, \(u = -n_j\) for one second and then \(u = 0\) keeps the
maximum critical position at \(\tfrac{34}{25} + \tfrac12 = \tfrac{93}{50} =
1.86 < 2\): all three physical states are indefinitely viable under full
information. Take three witness labels, each using its own position facet at
\(t_* = \tau + 1\), with kernels \(k_j(s) = \mathbf{1}_{\{s \le t_*\}} (t_* -
s) n_j\) and \(\beta_j = \tfrac{16}{25} - t_*\), weighted by
\(\lambda_1 = \tfrac38\), \(\lambda_2 = \lambda_3 = \tfrac{5}{16}\), so that
\(\sum_j \lambda_j n_j = 0\). Before observation the pooled kernel is
identically zero; after observation each mode has its own control and
\(\phi_U(\lambda_j (t_* - s) n_j) = -\lambda_j (t_* - s)\). Hence
\[
\Gamma(\tau) = -\int_\tau^{\tau+1} (\tau + 1 - s)\, ds
- \Bigl(\tfrac{16}{25} - \tau - 1\Bigr) = \tau - \frac{7}{50},
\tag{16}
\]
so \(\tau > \tfrac{7}{50}\) implies \(B_0 \notin \mathcal{K}_{\mathcal{I}}\).
Conversely, \(u = 0\) before observation and \(u = -n_j\) for one second
afterward keeps the maximum critical position at \(\tfrac{34}{25} + \tau +
\tfrac12\), safe exactly when \(\tau \le 0.14\). Thus \(\tau_{\max} =
\tfrac{7}{50} = 0.14\) is exact, and for horizons containing the stopping
time \(J_{\mathcal{I}}(\tau + 1) = \tau - 0.14\).

\subsubsection{6.2 Pairs are viable; the blind-window exit test misses}\label{pairs}

At \(\tau = \tfrac15\), the common control \(u = 0\) keeps all three branches
safe throughout the blind window: \(n_j^\top p_j(0.2) = \tfrac{39}{25} = 1.56
< 2\). No exit occurs before information arrives; nevertheless
\(\Gamma(0.2) = 0.06 > 0\) --- the best post-observation recourse cannot
recover all branches. The obstruction is not that information arrives after
unavoidable physical exit; it is that information arrives after some
indistinguishable branch has necessarily accumulated more braking liability
than the remaining control authority can discharge.

Every pair remains viable. For \(\{1,2\}\) use the common blind control
\(w = (-\tfrac25, -\tfrac45)\); reflect the second component for \(\{1,3\}\);
for \(\{2,3\}\) use \(w = (1,0)\); after revelation brake the normal velocity
with \(-n_j\), then the remaining tangential velocity. With \(\alpha = -n_j
^\top w\), the maximum critical position is \(\tfrac{34}{25} + \tau -
\alpha\tau^2/2 + (1 - \alpha\tau)^2/2\), giving \(1.9752\) for
\(\{1,2\},\{1,3\}\) (\(\alpha = 0.4\)) and \(1.9352\) for \(\{2,3\}\)
(\(\alpha = 0.6\)), both below \(2\); all values are exact rationals
(\(12345/6250\) and \(12193/6250\)). Over the seven nonempty priors formed
from the three states, exactly six are viable at delay \(\tfrac15\) --- a
calculation over a prior family, not a claim that all seven are reachable
beliefs from one fixed initial record.

\subsubsection{6.3 The finite program and its Farkas certificate}\label{finst}

Set \(\tau = \tfrac15\), \(T = \tfrac65\), \(h = \tfrac1{10}\): two common
pre-observation control blocks and ten post-observation blocks for each of
three modes, so \(Q = 32\), \(mQ = 64\). For the endpoint position rows the
exact moment error is \(e = 12 \cdot h^2/4 = \tfrac{3}{100}\) (the kernel is
affine on each of the \(T/h = 12\) mesh intervals and
\(h_U(\alpha n_j) = |\alpha|\)); the endpoint right-hand side is \(\beta =
\tfrac{16}{25} - \tfrac65 = -\tfrac{14}{25}\), relaxed to \(\beta + e =
-\tfrac{53}{100}\). The exact continuous certificate has margin
\(\tfrac{3}{50} = 0.06\); the finite moment certificate has margin
\(\tfrac{3}{100} = 0.03\).

The dual witness is exact. Order the input facets by
\(F = [(0,1);(0,-1);(2,1);(2,-1);(-2,1);(-2,-1)]\), \(f = (\tfrac45, \tfrac45,
2, 2, 2, 2)\), and define \(\eta_1 = (0,0,0,0,\tfrac14,\tfrac14)\),
\(\eta_2 = (0,\tfrac12,0,\tfrac3{10},0,0)\),
\(\eta_3 = (\tfrac12,0,\tfrac3{10},0,0,0)\), satisfying \(F^\top \eta_j =
-n_j\) and \(f^\top \eta_j = 1\). For a post-observation interval \(J_k\)
put \(a_k = \int_{J_k} (T - s)\, ds\) and \(\mu_{j,k} = \lambda_j a_k
\eta_j\). Before observation, stationarity holds because \(\sum_j \lambda_j
n_j = 0\); after observation, because \(\lambda_j a_k n_j + F^\top \mu_{j,k}
= 0\). Since \(\sum_k a_k = \tfrac12\), the Farkas contradiction is
\(\lambda^\top(\beta + e) + \mu^\top f_{\mathrm{blk}} = -\tfrac{53}{100} +
\tfrac12 = -\tfrac{3}{100} < 0\) --- an exact rational certificate. The
primal witness is equally exact: hold the pre-observation blocks at \(0\) and
the post-observation blocks of mode \(j\) at \(-n_j\); the endpoint row then
reads \(-\tfrac12 + \tfrac{14}{25} - \tfrac{3}{100} = \tfrac{3}{100}\), and
every other row lies at or below the same value (verified row by row in exact
arithmetic).

\subsubsection{6.4 Solver-certified mesh study}\label{mesh}

The joint audit of the predecessor material left one verification step
unexecuted: the mesh values below were known to match a closed-form lower
bound, but LP optimality had not been certified by an actual solver run. That
step is closed here. For each mesh the program (5) is built at its claimed
dimensions and solved by HiGHS; optimality is then certified in exact rational
arithmetic by the witness pair above (primal feasibility of every row;
dual stationarity and dual objective \(-\rho(h)\)), so the duality gap closes
exactly and the solver's returned value is a computational confirmation of a
proved optimum, not an observation. Table~\ref{tab:mesh} reports the
campaign: exact values \(\rho(h) = \max\{0,\ \tfrac{3}{50} - Th/4\}\),
solver values, iteration counts, and wall-clock times as measured on the
verification machine.

\begin{table*}[t]
\centering
\footnotesize
\begin{tabular}{@{}r|r|r|r|r|r|r|r@{}}
\toprule
\(h\) & blocks \(Q\) & variables \(mQ{+}1\) & rows \(Ms(N_t{+}1)+p_UQ\) & exact \(\rho(h)\) & solver \(\rho\) & iterations & time (s) \\
\midrule
0.20 & 16 & 33  & 243  & \(0\) & 0.000000000 & 36   & 0.024 \\
0.10 & 32 & 65  & 465  & \(\tfrac{3}{100} = 0.030\) & 0.030000000 & 78   & 0.014 \\
0.05 & 64 & 129 & 909  & \(\tfrac{9}{200} = 0.045\) & 0.045000000 & 141  & 0.043 \\
0.02 & 160 & 321 & 2241 & \(\tfrac{27}{500} = 0.054\) & 0.054000000 & 501  & 0.252 \\
0.01 & 320 & 641 & 4461 & \(\tfrac{57}{1000} = 0.057\) & 0.057000000 & 1062 & 1.015 \\
\bottomrule
\end{tabular}
\caption{Solver-certified mesh study on the three-branch instance
(\(\tau = \tfrac15\), \(T = \tfrac65\)): program dimensions, exact optimal
values, HiGHS-attained values (agreement to at most \(4.2 \times 10^{-17}\)),
simplex iterations, and wall-clock time. Optimality of each exact value is
certified by the exact primal--dual witness pair of Section~6.3, independent
of the solver. The exact continuous value is \(0.06\); three safety labels
support each positive certificate.}
\label{tab:mesh}
\end{table*}

The values obey the exact law \(\rho(h) = \tfrac{3}{50} - Th/4\) for
\(h < \tfrac15\) and vanish at \(h = \tfrac15\): the whole refinement gap is
the control-moment error \(Th/4\) of the endpoint rows. All dimension counts
(\(mQ + 1\) variables; \(Ms(N_t{+}1) + p_U Q\) inequalities), all witness
inequalities, and the closed-form law are machine-verified in exact rational
arithmetic by the companion verification script, which regenerates
Table~\ref{tab:mesh} verbatim and re-runs the solver campaign.

\subsection{7. Computation, complexity, and independent verification}\label{computation}

With \(Q\) information--time blocks and \(N_t + 1\) time labels, the program
has \(mQ + 1\) variables and at most \(Ms(N_t{+}1) + p_U Q\) inequalities
(\(p_U\) input-polytope facets); it is polynomial-time in these dimensions
and the rational data size, conditional on the validated coefficient bounds.
Initial-state dependence is affine, so no vertex enumeration of \(P_\theta\)
is required: at each label solve \(\max_{\xi \in P_\theta} c_i^\top
X_\theta(t,0)\xi\) and store a maximizing or near-maximizing point ---
the verifier needs only membership of the stored point in \(P_\theta\), not
global optimality; near-optimality bears on the convergence guarantee, not on
a positive certificate's validity. The same holds for stored disturbance
signals.

For rational outer rows \(\widehat A v \le \widehat b\),
\(F_{\mathrm{blk}} v \le f_{\mathrm{blk}}\), a residual-bounded certificate
test suffices: with \(\lambda, \mu \ge 0\), \(r = \widehat A^\top \lambda +
F_{\mathrm{blk}}^\top \mu\), \(c = \lambda^\top \widehat b + \mu^\top
f_{\mathrm{blk}}\), the condition
\[
c + \sum_j h_U(-r_j) < 0
\tag{15}
\]
is sufficient, since any feasible \(v\) satisfies \(-\sum_j h_U(-r_j) \le
r^\top v \le c\). All integration, rounding, and model-enclosure errors must
be included before applying the test, and uncertain row data must be relaxed
toward a \emph{larger} feasible controller set: overestimating the worst
disturbance contribution can create a false obstruction. A deletion-based
sparsification uses at most one feasibility solve per safety row; exact rank
compression requires verified rank, and floating-point singular-value
truncation must be charged to the error budget. If
\(J_{\mathcal{I}}(T) = \gamma > 0\), a mesh with \(R_U V_U h_u + R_D V_D h_d
+ L h_t < \gamma\) guarantees detection --- an explicit, though not claimed
optimal, margin-dependent certificate-size bound.

\subsection{8. Information--time rank: input dimension does not bound the witnesses}\label{rank}

\begin{proposition}[arbitrarily large minimal obstructions with one input]\label{prop:rank}
For every integer \(q \ge 1\) there is a continuous-time system with scalar
input \(u(t) \in [0,1]\) and \(q+1\) indistinguishable modes such that every
proper subbelief is viable, the full \((q{+}1)\)-mode belief is not, and every
obstruction certificate must involve all \(q+1\) modes.
\end{proposition}

Let \(T = q\), \(I_j = [j-1, j]\). For mode \(j \le q\):
\(\dot x_j = \mathbf{1}_{I_j}(t)(u - 1)\), \(x_j(0) = 0\); for mode \(q+1\):
\(\dot x_{q+1} = -u\), \(x_{q+1}(0) = q - \varepsilon\), \(0 < \varepsilon <
1\). Require \(0 \le x \le q\); there are no observations. Writing \(w_j =
\int_{I_j} u\), safety requires \(-w_j \le -1\) for \(j \le q\) and \(\sum_j
w_j \le q - \varepsilon\); summing gives \(0 \le -\varepsilon\), a
contradiction. If mode \(q+1\) is removed, \(u \equiv 1\) is safe; if mode
\(j \le q\) is removed, \(u = 0\) on \(I_j\) and \(u = 1\) elsewhere is safe.
Equal weights give certificate margin \(\varepsilon/(q+1)\); the kernel rank
is \(q\), and \(q+1\) witnesses are necessary.

A statement such as ``at most \(m+1\) compatible states suffice to witness
nonviability'' is valid for convex common-action sets in \(\mathbb{R}^m\) but
false for a common blind control function: the correct dimension counts
independent temporal and informational control decisions. This distinction is
consequential for review intervals, delayed sensing, and certificate-size
claims; its discrete counterpart is recorded as a scoping caveat in the
companion paper.

\subsection{9. Refinement erosion and oracle recourse}\label{erosion}

Suppose a coarse information cell \(C\) is split into finer cells
\(C_1,\dots,C_k\). Because
\(\phi_U(g_1 + \dots + g_k) \ge \sum_j \phi_U(g_j)\), the same witness
satisfies \(\Gamma_{\mathrm{coarse}} \ge \Gamma_{\mathrm{fine}}\), with the
exact difference
\[
\int \Bigl[ \phi_U\Bigl(\sum_j g_{C_j}\Bigr) - \sum_j \phi_U(g_{C_j}) \Bigr]
dt.
\tag{13}
\]
For \(U = [-\bar u, \bar u]^m\) this is \(\bar u \int \bigl[\sum_j
\|g_{C_j}(t)\|_1 - \|\sum_j g_{C_j}(t)\|_1\bigr] dt\): the certificate
identifies the precise time intervals on which indistinguishable branches
demand conflicting control effort. On rational instances the identity is
exactly checkable; the verification record confirms \(\phi_U(g_1 + g_2) \ge
\phi_U(g_1) + \phi_U(g_2)\) on \(3{,}240\) deterministic rational pairs for
the hexagonal \(U\) of Section~6.

The exact theorem assumes exogenous mode observations. A useful sound
extension does not: if finitely many admissible initial-state/disturbance
scenarios are provably indistinguishable until \(\tau\), give the controller
after \(\tau\) \emph{separate arbitrary control signals for every stored
scenario} --- an oracle relaxation of the true policy class. If
\[
\int_0^\tau \phi_U\Bigl(\sum_\ell \lambda_\ell k_\ell\Bigr) dt
+ \sum_\ell \int_\tau^T \phi_U(\lambda_\ell k_\ell)\, dt
- \sum_\ell \lambda_\ell \beta_\ell > 0,
\tag{14}
\]
with \(\beta_\ell\) computed from the explicitly stored scenario, then no
actual output-feedback policy can be safe. This covers arbitrary sensing
after a known blackout; it is sound but generally incomplete. The
three-branch instance realizing this recourse mode in the discrete setting is
in the companion paper.

\subsection{10. Scope delimitations}\label{scope}

The exact theorem does not cover: control-dependent observations or active
sensing; arbitrary continuous-state measurements; nonlinear dynamics without
additional validated enclosures; nonconvex actuator sets without a relaxation
gap; rate-constrained actuators without modifying the approximation; or
unrestricted infinite-horizon completeness. For general sensing after a
blackout, the oracle-recourse certificate is sound but incomplete. For
nonlinear systems, a uniform validated flow-error bound can be added to the
row budgets, yielding a sound extension, not automatic completeness. Against
the alternatives: full-information viability decides none of the
information-constrained question; initial admissibility and tangency tests
see no obstruction here; safety-before-first-observation tests are silent
when a common safe blind control exists; exact finite safety games await a
justified continuous-to-finite bridge, which is what Theorem~1 provides; and
na\"ive held-control infeasibility excludes only the held-control class.
No general-purpose computational-calculus claim is made: the theorem is a
certificate generator for the stated class, with completeness only under
refinement and semidecidability only with effectively computable bounds.

\subsection{References}\label{references}
{\footnotesize

Abaee, A.: An obstruction calculus for viability under incomplete observation.
Preprint. \url{https://github.com/MIKEAA2020/general-sustainability} (2026)

Anderson, E.J., Nash, P.: Linear Programming in Infinite-Dimensional Spaces:
Theory and Applications. Wiley, Chichester (1987)

Aubin, J.-P.: Viability Theory. Birkh\"auser, Boston (1991)

Cardaliaguet, P., Quincampoix, M., Saint-Pierre, P.: Set-valued numerical
methods for optimal control and differential games. In: Stochastic and
Differential Games: Theory and Numerical Methods, pp.~177--247. Birkh\"auser,
Boston (1999)

Hettich, R., Kortanek, K.O.: Semi-infinite programming: theory, methods, and
applications. SIAM Rev. \textbf{35}(3), 380--429 (1993)

Huangfu, Q., Hall, J.A.J.: Parallelizing the dual revised simplex method.
Math. Program. Comput. \textbf{10}(1), 119--142 (2018)

Veliov, V.M.: Sufficient conditions for viability under imperfect
measurement. Set-Valued Anal. \textbf{1}, 305--317 (1993)

}

\section*{Declarations}

\textbf{Funding.} None.\quad \textbf{Competing interests.} None.\quad
\textbf{Data availability.} No external data were used.\quad
\textbf{Code availability.} The verification scripts
\texttt{paper2c\_lp\_campaign.py} (solver-certified mesh campaign and exact
witnesses; regenerates Table~\ref{tab:mesh} verbatim) and
\texttt{paper2\_nonstandard\_verification.py} (51-check exact-arithmetic
record for the predecessor material) are publicly available at
\url{https://github.com/MIKEAA2020/general-sustainability} (folder
\texttt{arena agent 1/paper rewrites/latex}).\quad
\textbf{AI declaration.} GLM (Z.ai), Qwen (Alibaba Cloud) and DeepSeek AI
assisted with drafting and iterative review.

\end{document}
